% Ad Atticum 7.11 (ad-atticum-07-11) --- English translation
% Translated by Claude (Anthropic) from the Latin (Perseus).
% Style guide: STYLE.md.

\ciceroLetterOpener{CICERO ATTICO salutem}

\ciceroSection{1} I beg you, what is this? Or what is going on? For me it is all darkness. ``Cingulum,'' he says, ``we hold; Ancona we have lost; Labienus has left Caesar.'' Are we talking about a commander of the Roman people, or about Hannibal? Insensate, miserable man --- who has never so much as seen the shadow of \textit{[Greek: the honourable]}! And he says he does all this for the sake of his standing. But where is standing to be found, except where there is honour? Honourable, then, to keep an army with no public authority; to seize cities of citizens, so as to make easier the road to his fatherland; to engineer \textit{[Greek: cancellations of debt]}, \textit{[Greek: the recalls of exiles]}, and six hundred crimes besides, \textit{[Greek: that he may have the greatest of the gods, tyranny]}, for his own. Let him keep his fortune to himself! For my part, by Hercules, I would rather have one hour's basking with you in that little sun-trap of yours than all kingships of that sort --- or rather die a thousand deaths than once think anything of that kind. ``And what if you wanted it?'' you say.

\ciceroSection{2} Come now: who is there for whom wanting is not allowed? But this very wanting I think more miserable than being raised on a cross. There is one thing more miserable: to attain what you have so wanted. But enough of this. For amid these troubles I gladly \textit{[Greek: take so much leisure for it]}.

\ciceroSection{3} Let us go back to our man. By all that is holy! What kind of counsel of Pompey's does it seem to you --- this, I mean, that he has abandoned the City? For my part, \textit{[Greek: I am at a loss]}. Then nothing more absurd. You, abandon the City? Then the same, I suppose, if the Gauls were coming? ``The commonwealth,'' he says, ``does not lie in walls.'' But in altars and hearths it does. ``Themistocles did it.'' Yes --- because one city could not bear the wave of the whole barbarian world. But Pericles did not do the same, about fifty years later, when he held nothing but the walls; and our own forebears in the past, when the rest of the City was taken, kept the citadel all the same. \textit{[Greek: So it goes, somehow.]}

\ciceroSection{4} On the other hand, from the grief of the country-towns and from the talk of the men I meet, this plan looks as though it will turn out well. The complaint among people is extraordinary (I do not know whether it is so in Rome, but you will let me know) that the City is without magistrates, without a senate. Pompey, in flight, moves men wonderfully. What more can I say? The case has become a different one. They now think nothing should be conceded to Caesar. Explain to me what all this comes to.

\ciceroSection{5} I am in charge of a business that is not turbulent. For Pompey wants me to be the man whom this whole of Campania and the seaboard shall have as their \textit{[Greek: overseer]}, to whom the levy and the chief matters of business shall be referred. So I was thinking of being on the move. You can already see, I think, what the \textit{[Greek: onset]} of Caesar is, what the people is doing, what the state of the whole affair stands at. I should like you to write me of these things, and --- since they are changeable --- as often as you can. For I find rest both in writing to you and in reading what comes from you.
